\section{Before Jerome: The Plural Old Latin World}

Latin Christianity possessed Scripture before Jerome, but not as a single authorized Latin Bible. From at least the second century, biblical books circulated in Latin forms translated principally from Greek. Different translators, local revisions, oral and liturgical memory, and ordinary copying produced a field of related but non-identical texts. The modern label \emph{Vetus Latina} gathers that evidence for study; it must not conceal its plurality (Vetus Latina Institute; Houghton 2016, 7--30).

\subsection{Why translation began from Greek}

The earliest Latin Christian communities received the New Testament in Greek and commonly encountered the Old Testament through Greek Jewish Scriptures. A Latin Old Testament reading could therefore stand at two removes from a Hebrew form: a Latin translator rendered a particular Greek reading, and that Greek reading had its own textual and translational history. This is not evidence of carelessness. Greek Scripture was the Bible read, preached, and argued from in many ancient churches. A later comparison with Hebrew cannot simply be retrojected as the original operating rule of Latin Christianity.

The surviving material is uneven. Whole early Latin Bibles are absent; some books survive in fragmentary manuscripts or palimpsests, while patristic citations and liturgical texts preserve forms otherwise lost. A Father's wording may reflect memory, adaptation to an argument, or a manuscript no longer extant. It is still evidence, but evidence whose unit and use must be identified. The Beuron Vetus Latina project accordingly works book by book and witness by witness rather than printing a conjectural uniform ``Old Latin Bible.''

\evidenceframe{The Vetus Latina Institute's manuscript register and edition program.}{Pre-Jerome and non-Vulgate Latin evidence is dispersed among manuscripts, fragments, patristic citations, and liturgy; only part of the Bible has been critically edited in the series.}{The modern institutional category organizes surviving evidence. It does not prove that ancient readers recognized one edition bearing that name.}

\subsection{Variation as a late-antique problem}

Jerome's Gospel preface and Letter 27 describe numerous divergent Latin copies and defend correction from Greek manuscripts. His complaint is partisan testimony from a reviser under criticism, but it confirms that textual plurality was visible in Rome. Augustine, writing from North Africa, likewise treats comparison among Latin translations as useful and insists that faulty Latin copies be corrected by Greek, with the Septuagint important for the Old Testament and Greek manuscripts for the New (\emph{De doctrina christiana} 2.11--15).

Augustine's preference for the \emph{Itala} has sometimes been made the title of one identifiable Old Latin recension. The passage cannot securely bear that weight. It establishes his favorable judgment about a Latin form or quality described by that word; the long history of overlapping Latin texts makes a single recoverable edition uncertain. ``African,'' ``European,'' and similar classifications can be useful descriptions of affinities, but they are not self-contained Bibles whose boundaries are known in advance (Houghton 2016, 7--30).

\subsection{Books, not yet a stable pandect}

Biblical books and groups of books commonly circulated separately. Gospel books, Pauline collections, Psalters, and lectionary excerpts could develop different textual histories before later being gathered into pandects. A copyist might correct an Old Latin book toward a Vulgate form; a Vulgate copy could acquire a familiar Old Latin reading from worship or commentary. The direction of influence must therefore be demonstrated for each witness.

This material setting matters for the later word ``Vulgate.'' A reader who imagines one Old Latin Bible being replaced on a fixed date by one new Bible misses the mechanisms of change. Jerome could revise the Gospels without touching Acts. His Hebrew Psalter could circulate beside a Greek-based Psalter. A later complete Bible could combine parts of different origin. The codex eventually made such combinations visible as one physical book, but physical unity did not erase textual ancestry.

\subsection{The earlier meaning of ``vulgate''}

In Jerome's own usage, a ``common'' or ``vulgate'' edition could mean the widespread Septuagintal text, including its Old Latin descendants. The label had not yet become a proper name for his collected work. The reversal is historically revealing: the text that later generations called Jerome's Vulgate first challenged forms that a fourth-century writer could call vulgate. Later success renamed the field.

Two consequences govern the chapters that follow. First, Old Latin evidence is not merely debris swept away by a superior edition; it preserves early Latin reception and, at times, Greek readings of independent text-critical interest. Second, Jerome began inside this world. His revisions reused its wording, responded to its disputes, and depended on readers who already knew Latin Scripture by hearing and memory. Even a move toward Greek or Hebrew sources had to negotiate the authority of familiarity.
